\section{引言：从格式到内容设计}

随着 Agent 工具生态的快速发展，超过 30 种 Agent 工具（如 Claude Code、Gemini CLI、Cursor 等）已经在 \texttt{SKILL.md} 的文件格式上达成了标准化共识。格式问题基本已经解决，真正的挑战转移到了\textbf{内容设计}层面。

\begin{knowledgebox}{什么是 SKILL.md？}
\texttt{SKILL.md} 是 Agent Skills 规范中的核心文件，用于定义一个 skill 的指令、触发条件和行为逻辑。它是 Agent 工具加载和执行 skill 的入口。该规范由 Google ADK（Agent Development Kit）推动，已被多个主流 Agent 工具采纳。
\end{knowledgebox}

规范本身只说明了如何\textbf{打包}一个 skill，却没有提供如何\textbf{设计内部逻辑}的指导。例如，一个封装 FastAPI 约定的 skill 和一个四步文档生成 pipeline，虽然外部的 \texttt{SKILL.md} 文件看起来完全一样，但内部运作方式截然不同。

\begin{importantbox}{核心观点}
通过研究生态系统中 skill 的构建方式（从 Anthropic 的仓库到 Vercel 和 Google 的内部指南），可以总结出五种反复出现的设计模式，帮助开发者构建更可靠的 Agent。
\end{importantbox}

本文介绍的五种模式分别是：
\begin{enumerate}
    \item \textbf{Tool Wrapper}：让 Agent 成为任意库的即时专家
    \item \textbf{Generator}：基于可复用模板生成结构化文档
    \item \textbf{Reviewer}：按严重级别对代码进行检查评分
    \item \textbf{Inversion}：Agent 先采访你，再行动
    \item \textbf{Pipeline}：强制执行带检查点的多步骤工作流
\end{enumerate}

\subsection{本章小结}

Agent Skill 的格式标准化已经成熟，开发者的关注点应从"如何写对 YAML"转向"如何设计 skill 内部的逻辑结构"。五种设计模式为这一问题提供了系统性的解决框架。

